\chapter{Optimisation numérique en pratique}\label{ch:qm:numerical-optimisation-in-practice}

Chaque matin, un desk calibre un noyau à deux exponentielles sur une autocorrélation mesurée~: une composante rapide dont l'échelle de temps vaut environ trois jours
et une composante lente d'environ deux mois. Pendant la majeure partie de l'année, l'échelle de temps lente se situe entre 53 et 70 jours. Un lundi, elle affiche 88, alors que celle du vendredi
affichait 57, et l'ajustement aux données est, tout au plus, légèrement meilleur. Rien ne cloche dans l'optimiseur et rien ne s'est passé dans les données~: la
fonction objectif présente une longue vallée, presque plate, le long de laquelle l'échelle de temps lente et le poids de la composante rapide se compensent, et un peu de bruit a déplacé le
fond de la vallée. La même fonction objectif possède une vallée symétrique, atteinte en permutant les étiquettes des deux composantes, et des culs-de-sac au
bord du domaine. Ce chapitre traite des optimiseurs qui transforment de tels objectifs en nombres (méthodes de gradient et leur accélération, \omterm{def:qm:numerical-optimisation-in-practice:newton}{méthodes de Newton} et
quasi-Newton, \omterm{def:qm:numerical-optimisation-in-practice:lm}{moindres carrés non linéaires}, méthodes de décomposition et méthodes proximales, gradients stochastiques) et de la pratique qui rend une
\omterm{def:qm:numerical-optimisation-in-practice:calibration}{calibration} fiable~: transformations, \omterm{def:qm:numerical-optimisation-in-practice:calibration}{multistart}, diagnostics d'\omterm{def:qm:numerical-optimisation-in-practice:calibration}{identifiabilité} et pénalité ramenant vers la réponse de la veille.

\section{Méthodes du premier ordre}

\begin{definition}[Descente de gradient, recherche linéaire]\label{def:qm:numerical-optimisation-in-practice:gd}
La \emph{descente de gradient}\index{descente de gradient} itère $x_{k+1} = x_k - t_k\nabla f(x_k)$. Une \emph{recherche linéaire}\index{recherche linéaire} choisit le pas $t_k$ le long de la
direction de descente, par exemple pour satisfaire les conditions de Wolfe~: décroissance suffisante, $f(x + td) \le f(x) + c_1t\nabla f(x)^\top d$, et une condition de courbure,
$|\nabla f(x + td)^\top d| \le c_2|\nabla f(x)^\top d|$, avec $0 < c_1 < c_2 < 1$.
\end{definition}

\begin{proposition}[Vitesses de convergence de la descente de gradient]\label{prop:qm:numerical-optimisation-in-practice:rates}
Si $f$ est $\mu$-fortement convexe à gradient $L$-lipschitzien, la \omterm{def:qm:numerical-optimisation-in-practice:gd}{descente de gradient} à pas $1/L$ vérifie $f(x_k) - f^* \le (1 - \mu/L)^k(f(x_0) - f^*)$~: une convergence linéaire, à un taux
fixé par le conditionnement $\kappa = L/\mu$. Pour $f$ convexe sans forte convexité, l'écart décroît en $O(1/k)$.
\end{proposition}

\begin{proof}
Avec un pas $1/L$, le lemme de descente donne $f(x_{k+1}) \le f(x_k) - \frac1{2L}\lVert\nabla f(x_k)\rVert^2$, et la forte convexité donne $\lVert\nabla f(x_k)\rVert^2 \ge 2\mu(f(x_k) -
f^*)$~; en combinant, $f(x_{k+1}) - f^* \le (1 - \mu/L)(f(x_k) - f^*)$.
\end{proof}

\begin{definition}[Méthode du moment, accélération de Nesterov]\label{def:qm:numerical-optimisation-in-practice:momentum}
Une \emph{méthode du moment}\index{méthode du moment} ajoute au pas de gradient une fraction du pas précédent. L'\emph{accélération de Nesterov}\index{accélération de Nesterov} évalue le
gradient en un point extrapolé, $y_k = x_k + \frac{k-1}{k+2}(x_k - x_{k-1})$, $x_{k+1} = y_k - \nabla f(y_k)/L$, ce qui porte le taux convexe à $O(1/k^2)$ et le taux fortement convexe
à environ $1 - 1/\sqrt\kappa$ par itération.
\end{definition}

Sur la quadratique $\frac12(x_1^2 + 100x_2^2)$, la \omterm{def:qm:numerical-optimisation-in-practice:gd}{descente de gradient} simple depuis $(1, 1)$ avec un pas de $1/100$ élimine la coordonnée raide en une étape, puis réduit l'autre de $1 - 1/100$ par
pas, si bien que l'objectif passe de 50.5 à $\frac12(0.99)^{2k}$~: 0.067 après 100 itérations et 0.0090 après 200 (la borne de la proposition, $(1 - 1/\kappa)^k$, est pessimiste, à cause du carré)~;
la méthode de Nesterov atteint $5 \times 10^{-6}$ en 100 itérations (la \cref{fig:qm:numerical-optimisation-in-practice:rates}).
Le conditionnement est l'ennemi des méthodes du premier ordre, et le mauvais conditionnement est l'état normal des \omterm{def:qm:numerical-optimisation-in-practice:calibration}{calibrations} dont les paramètres se compensent.

\begin{omfigure}
\begin{tikzpicture}
\begin{semilogyaxis}[width=10cm,height=5cm,xlabel={itération},ylabel={$f(x_k) - f^*$},
  tick label style={font=\small},label style={font=\small},xmin=0,xmax=200,ymin=1e-12,ymax=10,grid=major,grid style={gray!25},
  legend columns=3,legend style={font=\footnotesize,at={(0.5,-0.3)},anchor=north,draw=none,fill=none,
  /tikz/every even column/.append style={column sep=8pt}},table/col sep=comma]
\addplot[omThm,very thick] table[x=iter,y=gd]{figdata/methods/24-numerical-optimisation-in-practice/rates.csv};
\addplot[black,dashed,thick] table[x=iter,y=theory]{figdata/methods/24-numerical-optimisation-in-practice/rates.csv};
\addplot[omDef,very thick] table[x=iter,y=nesterov]{figdata/methods/24-numerical-optimisation-in-practice/rates.csv};
\legend{descente de gradient,$\frac12(1 - 1/\kappa)^{2k}$,Nesterov}
\end{semilogyaxis}
\end{tikzpicture}

\omcaption{\omterm{def:qm:numerical-optimisation-in-practice:gd}{Descente de gradient} et méthode de Nesterov accélérée sur une quadratique de conditionnement $\kappa = 100$, pas $1/L$, depuis $(1, 1)$. La \omterm{def:qm:numerical-optimisation-in-practice:gd}{descente de gradient} suit exactement son taux linéaire~; la
méthode accélérée n'est pas monotone mais elle est beaucoup plus rapide. Données~: le tutoriel du chapitre.}\label{fig:qm:numerical-optimisation-in-practice:rates}
\end{omfigure}

\section{Newton, quasi-Newton et moindres carrés non linéaires}

\begin{definition}[Méthode de Newton, méthode de quasi-Newton, BFGS]\label{def:qm:numerical-optimisation-in-practice:newton}
La \emph{méthode de Newton}\index{méthode de Newton} se place directement au minimum du modèle quadratique local, $x_{k+1} = x_k - \nabla^2f(x_k)^{-1}\nabla f(x_k)$. Une \emph{méthode de quasi-Newton
méthode de quasi-Newton}\index{méthode de quasi-Newton} remplace la hessienne par une approximation mise à jour à partir des différences de gradients~; la \emph{méthode BFGS}\index{méthode BFGS} (Broyden, Fletcher, Goldfarb
et Shanno, 1970) met à jour l'approximation inverse $H$ par $H \leftarrow (I - \rho sy^\top)H(I - \rho ys^\top) + \rho ss^\top$, où $s$ désigne le pas, $y$ la variation du gradient et $\rho = 1/y^\top s$~;
L-BFGS (Liu et Nocedal, 1989) ne retient que les quelques dernières paires.
\end{definition}

Au voisinage d'un minimum dont la hessienne est définie positive, la \omterm{def:qm:numerical-optimisation-in-practice:newton}{méthode de Newton} converge quadratiquement~: le nombre de chiffres exacts double à chaque pas. La mise à jour BFGS garde $H$ définie positive
tant que $y^\top s > 0$, ce que garantit une \omterm{def:qm:numerical-optimisation-in-practice:gd}{recherche linéaire} satisfaisant la condition de courbure de Wolfe, puisqu'alors $y^\top s \ge (c_2 - 1)\nabla f(x)^\top s > 0$. Sur la quadratique précédente,
la \omterm{def:qm:numerical-optimisation-in-practice:newton}{méthode de Newton} ne fait qu'un pas et L-BFGS en fait sept.

\begin{definition}[Moindres carrés non linéaires, Gauss--Newton, Levenberg--Marquardt]\label{def:qm:numerical-optimisation-in-practice:lm}
Les \emph{moindres carrés non linéaires}\index{moindres carrés non linéaires} minimisent $\frac12\lVert r(x)\rVert^2$ pour un vecteur de résidus $r$. La \emph{méthode de Gauss--Newton}\index{méthode de Gauss--Newton}
approxime la hessienne par $J^\top J$, $J$ étant la jacobienne de $r$, et résout $J^\top J\,\Delta x = -J^\top r$. L'\emph{algorithme de Levenberg--Marquardt}\index{algorithme de Levenberg--Marquardt }
(Levenberg, 1944~; Marquardt, 1963) l'amortit, $(J^\top J + \lambda\operatorname{diag}(J^\top J))\Delta x = -J^\top r$, en augmentant $\lambda$ après un pas échoué (vers la \omterm{def:qm:numerical-optimisation-in-practice:gd}{descente de gradient}
mise à l'échelle) et en le diminuant après un succès (vers Gauss--Newton).
\end{definition}

La \omterm{def:qm:numerical-optimisation-in-practice:calibration}{calibration} par moindres carrés est le cas du chapitre~: les résidus sont l'autocorrélation du modèle moins celle mesurée aux retards 1 à 30. Le poids de la composante rapide appartient à $(0, 1)$
et les échelles de temps sont positives~: l'optimiseur travaille donc sur des transformations sans contrainte (logit et log) et ne propose jamais un paramètre inadmissible. $J^\top J$ à la solution mesure aussi
l'\omterm{def:qm:numerical-optimisation-in-practice:calibration}{identifiabilité}~: le premier jour, son conditionnement vaut 10.5 et la corrélation implicite entre le poids de la composante rapide et l'échelle de temps lente vaut 0.96, signature numérique de la
vallée plate.

\begin{definition}[Calibration, identifiabilité, multistart]\label{def:qm:numerical-optimisation-in-practice:calibration}
La \emph{calibration}\index{calibration} consiste à choisir les paramètres d'un modèle pour reproduire des quantités observées (prix, cotations, moments), en général par \omterm{def:qm:numerical-optimisation-in-practice:lm}{moindres carrés non linéaires}. Un paramètre est
localement non identifiable, et l'\emph{identifiabilité}\index{identifiabilité} fait défaut lorsque l'objectif est plat dans une direction à l'optimum, de sorte que les données ne peuvent pas le fixer. Le \emph{multistart}\index{multistart}
lance un optimiseur local à partir de nombreux points de départ pour explorer des minima distincts.
\end{definition}

Depuis 100 départs aléatoires sur la courbe du premier jour, 28 essais aboutissent au meilleur ajustement (erreur quadratique moyenne 0.0096), 35 aboutissent à son image miroir, avec les étiquettes des deux composantes permutées (la
même courbe, la même erreur, des paramètres différents), 26 divergent vers le bord du domaine, où l'échelle de temps lente atteint le cap de 1\,000 jours et où l'erreur d'ajustement est au moins deux fois plus grande, et 11 s'arrêtent
ailleurs (la \cref{fig:qm:numerical-optimisation-in-practice:multistart}). Une \omterm{def:qm:numerical-optimisation-in-practice:calibration}{calibration} menée depuis un seul point de départ fixe ne découvre jamais rien de tout cela.

\begin{omfigure}
\begin{tikzpicture}
\begin{loglogaxis}[width=9cm,height=5.2cm,xlabel={échelle de temps rapide (jours)},ylabel={échelle de temps lente (jours)},
  tick label style={font=\small},label style={font=\small},xmin=0.1,xmax=1500,ymin=0.1,ymax=1500,grid=major,grid style={gray!25},
  colorbar,colormap/viridis,point meta min=9,point meta max=60,colorbar style={ylabel={EQM ($\times 10^{-3}$)},ylabel style={font=\footnotesize},
  tick label style={font=\footnotesize}},table/col sep=comma]
\addplot[scatter,only marks,mark=*,mark size=2pt,scatter src=explicit] table[x=tau_fast,y=tau_slow,meta=rmse]{figdata/methods/24-numerical-optimisation-in-practice/multistart.csv};
\end{loglogaxis}
\end{tikzpicture}

\omcaption{Points d'arrivée de 100 essais de Levenberg--Marquardt à partir de départs aléatoires sur l'autocorrélation d'un jour, tracés selon leurs deux échelles de temps (triées, de sorte que les solutions symétriques coïncident) et colorés selon
l'erreur d'ajustement~: un amas au meilleur ajustement (3 et 55 jours~; 63 essais, images miroir comprises), des essais bloqués sur le cap de 1\,000 jours, et quelques-uns ailleurs. Données~: le tutoriel du chapitre,
graine fixée.}\label{fig:qm:numerical-optimisation-in-practice:multistart}
\end{omfigure}

La vallée elle-même apparaît dans un profil (la \cref{fig:qm:numerical-optimisation-in-practice:valley})~: en fixant l'échelle de temps lente et en réoptimisant les deux autres paramètres, l'erreur d'ajustement vaut
$9.9 \times 10^{-3}$ à 62 jours, $12.3 \times 10^{-3}$ à 42 et $12.9 \times 10^{-3}$ à 89, contre un niveau de bruit de $10 \times 10^{-3}$. Toute valeur comprise entre 45 et 80 jours environ reste dans le niveau de bruit.

\begin{omfigure}
\begin{tikzpicture}
\begin{semilogxaxis}[width=10cm,height=5cm,xlabel={échelle de temps lente $\tau_2$ (jours), fixée},ylabel={meilleur RMSE ($\times 10^{-3}$)},
  tick label style={font=\small},label style={font=\small},xmin=20,xmax=400,ymin=5,ymax=25,grid=major,grid style={gray!25},
  xtick={20,50,100,200,400},xticklabels={20,50,100,200,400},table/col sep=comma]
\addplot[omDef,very thick] table[x=tau2,y=rmse]{figdata/methods/24-numerical-optimisation-in-practice/valley.csv};
\addplot[black,dashed,thin,sharp plot] coordinates {(20,10) (400,10)};
\end{semilogxaxis}
\end{tikzpicture}

\omcaption{Profil de l'objectif de \omterm{def:qm:numerical-optimisation-in-practice:calibration}{calibration} le long de l'échelle de temps lente~: pour chaque $\tau_2$ fixé, la meilleure erreur quadratique moyenne racine sur les deux autres paramètres, sur la courbe d'un jour. La ligne pointillée
est le niveau de bruit de l'autocorrélation mesurée. Données~: le tutoriel du chapitre, graine fixée.}\label{fig:qm:numerical-optimisation-in-practice:valley}
\end{omfigure}

Recalibrer chaque jour sur 250 jours à partir du même départ~: l'échelle de temps lente présente une variation d'un jour à l'autre de médiane 4.5 jours et de maximum 30.7 (jour 189~: de 56.8 à 87.6 jours, erreur d'ajustement de
0.00926 à 0.00921). Avec une pénalité de Tikhonov vers les paramètres de la veille (poids 0.01 sur le carré de la variation de chaque paramètre transformé, l'ajustement partant de la réponse de la veille), la
variation médiane vaut 1.9 jour et le maximum 9.5, l'écart type de l'échelle de temps lente sur l'année passe de 5.5 à 2.7 jours, et l'erreur d'ajustement moyenne passe de 0.00945 à 0.00952, soit une hausse de
0.7~\% (la \cref{fig:qm:numerical-optimisation-in-practice:daily}).

\begin{omfigure}
\begin{tikzpicture}
\begin{axis}[width=12cm,height=5cm,xlabel={jour},ylabel={échelle de temps lente $\tau_2$ (jours)},
  tick label style={font=\small},label style={font=\small},xmin=1,xmax=250,ymin=40,ymax=100,grid=major,grid style={gray!25},
  legend columns=2,legend style={font=\footnotesize,at={(0.5,-0.3)},anchor=north,draw=none,fill=none,
  /tikz/every even column/.append style={column sep=8pt}},table/col sep=comma]
\addplot[omThm,thin] table[x=day,y=free]{figdata/methods/24-numerical-optimisation-in-practice/daily.csv};
\addplot[omDef,very thick] table[x=day,y=penalised]{figdata/methods/24-numerical-optimisation-in-practice/daily.csv};
\addplot[black,dashed,thin,sharp plot] coordinates {(1,60) (250,60)};
\legend{recalibration depuis un départ fixe,avec une pénalité vers la veille}
\end{axis}
\end{tikzpicture}

\omcaption{L'échelle de temps lente calibrée sur 250 jours de mesures bruitées du même noyau vrai (valeur vraie 60 jours, pointillés)~: recalibration chaque jour depuis un départ fixe, et recalibration à partir des
paramètres de la veille avec une pénalité de Tikhonov vers ceux-ci. Données~: le tutoriel du chapitre, graine fixée.}\label{fig:qm:numerical-optimisation-in-practice:daily}
\end{omfigure}

\section{Séparation d'opérateurs}

\begin{definition}[Opérateur proximal, méthode du gradient proximal]\label{def:qm:numerical-optimisation-in-practice:prox}
L'\emph{opérateur proximal}\index{opérateur proximal} d'une \omterm{def:qm:convex-optimisation:convex}{fonction convexe} $g$ est $\operatorname{prox}_{tg}(v) = \arg\min_x\bigl(g(x) + \frac1{2t}\lVert x - v\rVert^2\bigr)$~; pour $g = \lambda\lVert\cdot\rVert_1$, il s'agit d'un
seuillage doux. La \emph{méthode du gradient proximal}\index{méthode du gradient proximal} minimise $f + g$, $f$ étant lisse, par $x_{k+1} = \operatorname{prox}_{g/L}(x_k - \nabla f(x_k)/L)$~; FISTA (Beck et
Teboulle, 2009) ajoute l'extrapolation de Nesterov.
\end{definition}

Le \omterm{def:qm:linear-models-under-stress:penalties}{lasso} du chapitre~16 en est le cas d'école. Sur une régression de 200 observations portant sur 50 prédicteurs, dont deux presque colinéaires, la \omterm{def:qm:numerical-optimisation-in-practice:prox}{méthode du gradient proximal} ordinaire (ISTA) se trouve encore à
$7.6 \times 10^{-4}$ de l'objectif optimal après 50 itérations et à $4.6 \times 10^{-4}$ après 100~; FISTA n'en est plus qu'à environ $10^{-7}$ après 50 et à moins de $10^{-10}$ après 65
(la \cref{fig:qm:numerical-optimisation-in-practice:lasso}).

\begin{omfigure}
\begin{tikzpicture}
\begin{semilogyaxis}[width=10cm,height=5cm,xlabel={itération},ylabel={écart à l'objectif},
  tick label style={font=\small},label style={font=\small},xmin=0,xmax=150,ymin=1e-10,ymax=1,grid=major,grid style={gray!25},
  legend columns=2,legend style={font=\footnotesize,at={(0.5,-0.3)},anchor=north,draw=none,fill=none,
  /tikz/every even column/.append style={column sep=8pt}},table/col sep=comma]
\addplot[omThm,very thick] table[x=iter,y=ista]{figdata/methods/24-numerical-optimisation-in-practice/lasso.csv};
\addplot[omDef,very thick] table[x=iter,y=fista]{figdata/methods/24-numerical-optimisation-in-practice/lasso.csv};
\legend{gradient proximal (ISTA),accéléré (FISTA)}
\end{semilogyaxis}
\end{tikzpicture}

\omcaption{Le \omterm{def:qm:linear-models-under-stress:penalties}{lasso} ($\lambda = 0.1$, 200 observations, 50 prédicteurs dont une paire presque colinéaire) par la \omterm{def:qm:numerical-optimisation-in-practice:prox}{méthode du gradient proximal} et sa version accélérée~: écart à l'objectif optimal
en fonction de l'itération, plafonné à $10^{-10}$. Données~: le tutoriel du chapitre, graine fixée.}\label{fig:qm:numerical-optimisation-in-practice:lasso}
\end{omfigure}

\begin{definition}[Méthode des directions alternées des multiplicateurs]\label{def:qm:numerical-optimisation-in-practice:admm}
La \emph{méthode des directions alternées des multiplicateurs}\index{méthode des directions alternées des multiplicateurs} (ADMM) minimise $f(x) + g(z)$ sous la contrainte $x = z$ en alternant $x \leftarrow
\operatorname{prox}_{f/\rho}(z - u)$, $z \leftarrow \operatorname{prox}_{g/\rho}(x + u)$, $u \leftarrow u + x - z$ (Boyd, Parikh, Chu, Peleato et Eckstein, 2011).
\end{definition}

L'ADMM converge pour des fonctions $f$ et $g$ convexes fermées dès qu'une solution existe, à une vitesse modeste, et chaque étape ne coûte guère plus que les deux opérateurs proximaux. Pour la \omterm{def:qm:convex-optimisation:ncm}{matrice de corrélation la plus proche} présentée au
chapitre~23, $f$ est la distance à la matrice donnée augmentée de la contrainte de semi-définie positivité (un \omterm{def:qm:covariance-estimation-and-random-matrices:clip}{écrêtage des valeurs propres}) et $g$ la contrainte de diagonale unitaire. Sur la matrice défectueuse du chapitre~23, l'ADMM nécessite 91
itérations contre 74 pour les projections alternées, et les deux réponses coïncident à $1.4 \times 10^{-9}$ près.

\section{Gradients stochastiques}

\begin{definition}[Descente de gradient stochastique]\label{def:qm:numerical-optimisation-in-practice:sgd}
La \emph{descente de gradient stochastique}\index{descente de gradient stochastique} minimise $\frac1n\sum_if_i(x)$ par des pas le long du gradient d'un seul terme tiré au hasard (ou d'un petit lot),
$x_{k+1} = x_k - t_k\nabla f_{i_k}(x_k)$.
\end{definition}

Chaque pas est une estimation sans biais, mais bruitée, d'un pas de gradient. Avec un pas constant, les itérés oscillent dans une boule de bruit autour du minimum~; Robbins et Monro (1951) ont montré que des pas tels que
$\sum_kt_k = \infty$ et $\sum_kt_k^2 < \infty$, comme $t_k = a/(b + k)$, convergent. Sur un problème de moindres carrés à 1\,000 observations, 20\,000 pas aboutissent à 0.009 de la solution des moindres carrés
avec des pas de Robbins--Monro, et à 0.113 avec un pas constant de 0.05. Les méthodes adaptatives comme Adam (Kingma et Ba, 2014) recalibrent le pas de chaque coordonnée à partir des moments glissants de ses
gradients~; elles constituent le choix par défaut pour les réseaux de neurones du livre~9, et un mauvais choix par défaut pour une \omterm{def:qm:numerical-optimisation-in-practice:calibration}{calibration} à trois paramètres, où Levenberg--Marquardt converge en quelques dizaines de pas.

\section{Tutoriel~: le lundi où l'ajustement a bougé}\label{tut:qm:numerical-optimisation-in-practice:monday}

\begin{tutorial}
\textbf{Objectif.} Calibrer chaque jour un noyau à deux exponentielles, explorer ses minima par \omterm{def:qm:numerical-optimisation-in-practice:calibration}{multistart}, diagnostiquer la vallée plate et stabiliser l'ajustement quotidien par une pénalité vers la veille.
\textbf{État final~:} les \cref{fig:qm:numerical-optimisation-in-practice:multistart,fig:qm:numerical-optimisation-in-practice:valley,fig:qm:numerical-optimisation-in-practice:daily}~; le saut médian de
4.5 et 1.9 jours et le coût d'ajustement de 0.7~\%.
\begin{tutsteps}
\item \textbf{Levenberg--Marquardt} avec son calendrier d'amortissement.
\omcode{code/firm/optim/firm_optim.py}{117}{142}{Levenberg--Marquardt pour les moindres carrés non linéaires.}\label{lst:qm:numerical-optimisation-in-practice:lm}
\item \textbf{La pénalité vers la veille} et le diagnostic d'\omterm{def:qm:numerical-optimisation-in-practice:calibration}{identifiabilité}.
\omcode{code/firm/optim/firm_optim.py}{222}{239}{Diagnostics d'identifiabilité et pénalité de Tikhonov.}\label{lst:qm:numerical-optimisation-in-practice:tikhonov}
\item \textbf{Exécuter} \texttt{daily(0.0)}, \texttt{daily(0.01)}, \texttt{multistart()}, \texttt{valley()} et \texttt{diagnostics()} dans \texttt{qm\_optim.py}, puis \texttt{fig\_optim.py}.
\end{tutsteps}
\textbf{Pour aller plus loin.} Ordonner les échelles de temps par construction ($\tau_2 = \tau_1 + e^u$) pour supprimer la vallée miroir~; choisir le poids de la pénalité par \omterm{def:qm:linear-models-under-stress:cv}{validation croisée} sur l'erreur d'ajustement du lendemain.
\end{tutorial}

\section{Construction~: le harnais de calibration}\label{bld:qm:numerical-optimisation-in-practice:optim}

\begin{build}
\bfield{Rôle} Tout modèle que le petit cabinet calibre chaque jour (noyaux, surfaces de volatilité, courbes de taux) passe par un même harnais qui transforme, lance le \omterm{def:qm:numerical-optimisation-in-practice:calibration}{multistart}, pénalise et journalise.
\bfield{Interface} \texttt{gradient\_descent}, \texttt{newton}, \texttt{lbfgs}, \texttt{levenberg\_marquardt}, \texttt{prox\_l1}, \texttt{proximal\_gradient}, \texttt{admm\_nearest\_correlation},
\texttt{sgd}, \texttt{Bounded}, \texttt{multistart}, \texttt{identifiability}, \texttt{tikhonov}.
\bfield{Règles} Les paramètres sont optimisés dans des coordonnées non contraintes~; chaque \omterm{def:qm:numerical-optimisation-in-practice:calibration}{calibration} de production journalise son point de départ, ses itérations, son erreur d'ajustement, le conditionnement de $J^\top J$ et la variation depuis
la veille~; un nouveau minimum éloigné de celui de la veille est signalé, et non publié silencieusement~; le \omterm{def:qm:numerical-optimisation-in-practice:calibration}{multistart} tourne chaque semaine sur chaque modèle.
\bfield{Tests d'acceptation} \texttt{code/firm/optim/tests/}~: le taux linéaire de la \omterm{def:qm:numerical-optimisation-in-practice:gd}{descente de gradient}, l'accélération, Newton sur Rosenbrock~; L-BFGS sur Rosenbrock~; Levenberg--Marquardt retrouve un ajustement exponentiel
et une pénalité forte le fige~; seuillage doux et \omterm{def:qm:linear-models-under-stress:penalties}{lasso} par FISTA~; \omterm{def:qm:convex-optimisation:ncm}{matrice de corrélation la plus proche} par ADMM~; SGD avec les pas de Robbins--Monro~; les transformations font l'aller-retour~; ordre du \omterm{def:qm:numerical-optimisation-in-practice:calibration}{multistart}.
\bfield{Pour aller plus loin} Contraintes de boîte par projection dans L-BFGS (L-BFGS-B)~; différentiation automatique des résidus (chapitre~28)~; un Levenberg--Marquardt à région de confiance.
\end{build}

\begin{omsources}
\item K.~Levenberg, «~A method for the solution of certain non-linear problems in least squares~», \emph{Quarterly of Applied Mathematics} 2, 1944; D.~W.~Marquardt, «~An algorithm for least-squares
estimation of nonlinear parameters~», \emph{Journal of the Society for Industrial and Applied Mathematics} 11, 1963.
\item D.~C.~Liu et J.~Nocedal, «~On the limited memory BFGS method for large scale optimization~», \emph{Mathematical Programming} 45, 1989.
\item H.~Robbins et S.~Monro, «~A stochastic approximation method~», \emph{Annals of Mathematical Statistics} 22, 1951.
\item A.~Beck et M.~Teboulle, «~A fast iterative shrinkage-thresholding algorithm for linear inverse problems~», \emph{SIAM Journal on Imaging Sciences} 2, 2009.
\item S.~Boyd, N.~Parikh, E.~Chu, B.~Peleato et J.~Eckstein, «~Distributed optimization and statistical learning via the alternating direction method of multipliers~», \emph{Foundations and Trends in
Machine Learning} 3, 2011.
\item D.~P.~Kingma et J.~Ba, «~Adam: a method for stochastic optimization~», arXiv:1412.6980, 2014.
\item J.~Nocedal et S.~J.~Wright, \emph{Numerical Optimization}, 2\up{e} éd., Springer, 2006.
\end{omsources}

\section{Exercices}

\begin{exercise}[$\star$]\label{exo:qm:numerical-optimisation-in-practice:1}
Combien d'itérations de \omterm{def:qm:numerical-optimisation-in-practice:gd}{descente de gradient} réduisent-elles l'erreur d'une fonction quadratique de conditionnement 1\,000 d'un facteur $10^{6}$~? Et avec l'\omterm{def:qm:numerical-optimisation-in-practice:momentum}{accélération de Nesterov}, à peu près~?
\end{exercise}

\begin{exercise}[$\star$]\label{exo:qm:numerical-optimisation-in-practice:2}
Calculer $\operatorname{prox}_{t\lambda|\cdot|}(v)$ pour $v = (3, -0.5, 1)$, $t\lambda = 1$.
\end{exercise}

\begin{exercise}[$\star$]\label{exo:qm:numerical-optimisation-in-practice:3}
Pourquoi le noyau à deux exponentielles présente-t-il deux minima donnant exactement le même ajustement~?
\end{exercise}

\begin{exercise}[$\star\star$]\label{exo:qm:numerical-optimisation-in-practice:4}
Montrer qu'un pas de Newton depuis un point quelconque minimise une fonction quadratique strictement convexe.
\end{exercise}

\begin{exercise}[$\star\star$]\label{exo:qm:numerical-optimisation-in-practice:5}
Montrer que la condition de courbure de Wolfe implique $y^\top s > 0$ pour une direction de descente.
\end{exercise}

\begin{exercise}[$\star\star$]\label{exo:qm:numerical-optimisation-in-practice:6}
Les pas $t_k = 1/\sqrt k$ vérifient-ils les conditions de Robbins--Monro~? Et $t_k = 1/k$~?
\end{exercise}

\begin{exercise}[$\star\star\star$]\label{exo:qm:numerical-optimisation-in-practice:7}
\emph{Programmation.} Lancer la \omterm{def:qm:numerical-optimisation-in-practice:calibration}{calibration} quotidienne avec des poids de pénalité 0, 0.001, 0.01 et 0.1, et dresser un tableau de la variation quotidienne médiane de l'échelle de temps lente en regard de l'erreur d'ajustement moyenne.
\end{exercise}

\begin{exercise}[$\star\star\star$]\label{exo:qm:numerical-optimisation-in-practice:8}
\emph{Repérer le défaut.} «~Notre \omterm{def:qm:numerical-optimisation-in-practice:calibration}{calibration} converge toujours, avec une norme du gradient inférieure à $10^{-8}$, donc ses paramètres sont fiables.~»
\end{exercise}

\section{Problème~: le lundi où l'ajustement a bougé}

\begin{problem}[{Problème du week-end --- la vallée plate d'une calibration}]\label{pb:qm:numerical-optimisation-in-practice:1}
Chaque jour, un desk ajuste $\rho(k) = a e^{-k/\tau_1} + (1 - a)e^{-k/\tau_2}$, $k = 1, \dots, 30$, à une autocorrélation mesurée dont les vrais paramètres sont $a = 0.6$, $\tau_1 = 3$ et $\tau_2 = 60$ jours, mesurée
avec un bruit d'écart-type 0.01. L'ajustement se fait par Levenberg--Marquardt en coordonnées logit et logarithmiques. Deux cent cinquante jours sont simulés.

\medskip\noindent\textbf{Partie I --- Un jour.}
\begin{enumerate}
\item Depuis le point de départ fixe, que donne l'ajustement du premier jour, et avec quelle erreur~?
\item Que trouvent 100 points de départ aléatoires~?
\item Quelle est la solution symétrique, et pourquoi existe-t-elle~?
\item Quels sont le conditionnement de $J^\top J$ et la corrélation des paramètres~?
\item Sur quelle plage de l'échelle de temps lente l'ajustement reste-t-il dans le bruit~?
\end{enumerate}

\medskip\noindent\textbf{Partie II --- Chaque jour.}
\begin{enumerate}[resume]
\item Quelles sont les variations quotidiennes médiane et maximale de l'échelle de temps lente~?
\item Que s'est-il passé le jour de la variation maximale~?
\item Avec une pénalité vers la veille de poids 0.01, quelles sont-elles~?
\item Quel est le coût de la pénalité en erreur d'ajustement~?
\item De combien diminue l'étendue annuelle de l'échelle de temps lente~?
\end{enumerate}

\medskip\noindent\textbf{Partie III --- Méthodes.}
\begin{enumerate}[resume]
\item Comment se comparent la \omterm{def:qm:numerical-optimisation-in-practice:gd}{descente de gradient} et la méthode de Nesterov sur une fonction quadratique de conditionnement 100~?
\item Combien d'itérations la \omterm{def:qm:numerical-optimisation-in-practice:newton}{méthode de Newton} et L-BFGS demandent-elles sur ce problème~?
\item Comment se comparent ISTA et FISTA sur le \omterm{def:qm:linear-models-under-stress:penalties}{lasso}~?
\item Comment se comparent ADMM et les projections alternées pour la \omterm{def:qm:convex-optimisation:ncm}{matrice de corrélation la plus proche}~?
\item Qu'apportent les pas de Robbins--Monro par rapport à un pas constant~?
\end{enumerate}

\medskip\noindent\textbf{Partie IV --- Jugement.}
\begin{enumerate}[resume]
\item Quel paramètre faudrait-il publier avec une mise en garde, et laquelle~?
\item Comment choisir le poids de la pénalité~?
\item Comment supprimer la vallée symétrique~?
\item Énoncer le \emph{résultat nommé}~: le saut quotidien des paramètres avec et sans pénalité, et l'erreur d'ajustement que coûte la pénalité.
\item En une phrase~: que ne révèle pas une \omterm{def:qm:numerical-optimisation-in-practice:calibration}{calibration} qui a convergé~?
\end{enumerate}
\end{problem}

\section{Questions d'entretien}

\begin{interviewq}[$\star$ \iqroles{chercheur, ingénieur ML}]\label{iq:qm:numerical-optimisation-in-practice:1}
Pourquoi la \omterm{def:qm:numerical-optimisation-in-practice:gd}{descente de gradient} ralentit-elle sur des problèmes mal conditionnés, et que faites-vous pour y remédier~?
\end{interviewq}

\begin{interviewq}[$\star\star$ \iqroles{chercheur, développeur}]\label{iq:qm:numerical-optimisation-in-practice:2}
Les paramètres de votre \omterm{def:qm:numerical-optimisation-in-practice:calibration}{calibration} de modèle sautent d'un jour à l'autre alors que l'erreur d'ajustement reste la même. Que se passe-t-il, et que faites-vous~?
\end{interviewq}

\begin{interviewq}[$\star\star$ \iqroles{chercheur}]\label{iq:qm:numerical-optimisation-in-practice:3}
Expliquez l'\omterm{def:qm:numerical-optimisation-in-practice:lm}{algorithme de Levenberg--Marquardt}. Pourquoi ne pas utiliser simplement Gauss--Newton~?
\end{interviewq}

\begin{interviewq}[$\star\star$ \iqroles{ingénieur ML}]\label{iq:qm:numerical-optimisation-in-practice:4}
Pourquoi la \omterm{def:qm:numerical-optimisation-in-practice:sgd}{descente de gradient stochastique} à taux d'apprentissage constant ne converge-t-elle pas, et quelles sont les conditions de Robbins--Monro~?
\end{interviewq}

\begin{interviewq}[$\star\star$ \iqroles{chercheur, ingénieur ML}]\label{iq:qm:numerical-optimisation-in-practice:5}
Qu'est-ce qu'un \omterm{def:qm:numerical-optimisation-in-practice:prox}{opérateur proximal}, et pourquoi est-il utile pour le \omterm{def:qm:linear-models-under-stress:penalties}{lasso}~?
\end{interviewq}

\begin{interviewq}[$\star\star\star$ \iqroles{chercheur}]\label{iq:qm:numerical-optimisation-in-practice:6}
Pourquoi BFGS a-t-il besoin d'une \omterm{def:qm:numerical-optimisation-in-practice:gd}{recherche linéaire} satisfaisant les conditions de Wolfe~?
\end{interviewq}
